\documentclass[11pt]{article}
\usepackage[margin=1in]{geometry}
\usepackage{amsmath,amssymb,amsthm}
\usepackage{xurl}
\usepackage{hyperref}
\sloppy
\let\oldus\_ \renewcommand{\_}{\oldus\allowbreak}
\newtheorem{theorem}{Theorem}
\newtheorem{lemma}[theorem]{Lemma}
\theoremstyle{definition}
\newtheorem{definition}[theorem]{Definition}
\newcommand{\qb}[2]{\genfrac{[}{]}{0pt}{}{#1}{#2}_{q}}
\newcommand{\F}{\mathbb{F}}

\title{Proof of conjecture 00000002490:\\ subspace counts over finite fields are Gaussian binomials,\\ and so are the counts in every interval of the subspace lattice}
\author{Jackmeson1}
\date{2026-10-05}

\begin{document}
\maketitle

\section{The conjecture and its reading}

\textbf{Statement (English, verbatim).} Definition: The finite-field side of partitions:
q-binomial lattices. Conjecture: The counting of subspaces over finite fields is governed by
Gaussian binomial coefficients: the lattice count closes under the Gaussian binomial structure.
(partition finite-field side q-binomial lattice Gaussian binomials)

\textbf{Chinese text.} It says, in substance, that on the finite-field side of partitions the
counts of the lattice are Gaussian binomials.

\textbf{This is a classical theorem, not a new result.} We prove it and check the proof in
Lean 4 with Mathlib. The wording is vague, so we prove the following explicit statements. Each
one is a theorem in the Lean file.
\begin{itemize}
\item[(R1)] \emph{Counting.} Let $\F$ be any finite field with $q$ elements and let $n,k\ge 0$.
  The number of $k$-dimensional subspaces of $\F^n$ is $\qb{n}{k}$. More generally, the same
  holds for every $n$-dimensional $\F$-vector space $V$.
\item[(R2)] \emph{The counts close under the q-Pascal structure.} Write $N_q(n,k)$ for the
  number in (R1). Then both q-Pascal rules hold for these counts:
  $N_q(n+1,k+1)=N_q(n,k)+q^{k+1}N_q(n,k+1)$, and, for $k\le n$,
  $N_q(n+1,k+1)=N_q(n,k+1)+q^{n-k}N_q(n,k)$.
\item[(R3)] \emph{Interval closure.} Every interval of the subspace lattice is again counted by
  Gaussian binomials. Let $U\le T$ be subspaces of a finite $\F$-vector space $V$, and let
  $\dim U\le k$. Then the number of $k$-dimensional $W$ with $U\le W\le T$ is
  $\qb{\dim T-\dim U}{k-\dim U}$.
\end{itemize}
\textbf{Not formalized.} The word ``partitions'' appears only in the heading. We do not
formalize the separate fact that the coefficients of $\qb{n}{k}$ count partitions in a box.
``Lattice'' is read as the lattice of subspaces of $\F^n$, ordered by inclusion.

\textbf{Conventions.} $\F$ is an arbitrary finite field with $q=|\F|$ elements, so $q\ge 2$.
Its prime-power form is never used. A subspace is a Mathlib \texttt{Submodule}. Dimension is
\texttt{Module.finrank}. A count is \texttt{Nat.card} of a subtype. All $k\ge 0$ are allowed,
and for $k>n$ both sides are $0$. Natural-number subtraction is truncated, so $n-k$ and
$k-\dim U$ appear only under the hypotheses $k\le n$ and $\dim U\le k$. We write
$\binom{k}{2}=k(k-1)/2$.

\section{Definitions as formalized}

\begin{definition}[\texttt{gaussBinom}]
The Gaussian binomial $\qb{n}{k}\in\mathbb{N}[q]$ is defined by the rules
$\qb{n}{0}=1$, $\qb{0}{k+1}=0$, and
$\qb{n+1}{k+1}=\qb{n}{k}+q^{k+1}\qb{n}{k+1}$.
\end{definition}

Wikipedia [W] gives this rule as the first ``analog of Pascal's identity''. It writes the rule
as $\binom{m}{r}_q=q^r\binom{m-1}{r}_q+\binom{m-1}{r-1}_q$. Wikipedia \emph{defines} the
coefficient by the product formula
\[\binom{m}{r}_q=\frac{(1-q^m)(1-q^{m-1})\cdots(1-q^{m-r+1})}{(1-q)(1-q^2)\cdots(1-q^r)}.\]
We prove that our definition agrees with that one. Lemma~\ref{lem:prod} shows that
$\qb{n}{k}\cdot\prod_{i<k}(1-q^{i+1})=\prod_{i<k}(1-q^{n-i})$ holds in every commutative ring.
In $\mathbb{Z}[q]$ the denominator is nonzero, because its value at $q=0$ is $1$
(\texttt{qDen\_X\_ne\_zero}). So the product formula determines $\qb{n}{k}$ uniquely.

\section{Proof}

Let $P_x(n,k)=\prod_{i<k}(x^n-x^i)$ (Lean \texttt{liProd}).

\begin{lemma}[\texttt{gaussBinom\_mul\_liProd}]\label{lem:li}
In every commutative ring, $\qb{n}{k}(x)\,P_x(k,k)=P_x(n,k)$.
\end{lemma}
\begin{proof}
The proof is by induction on $n$, for all $k$ at once. The base cases are immediate, since
$P_x(0,k+1)$ has the factor $x^0-x^0=0$. The induction uses
$P_x(n+1,k+1)=(x^{n+1}-1)x^kP_x(n,k)$ and $P_x(n,k+1)=P_x(n,k)(x^n-x^k)$. With these, the
q-Pascal rule reduces the step to the identity
$(x^{k+1}-1)x^k+x^{k+1}(x^n-x^k)=x^k(x^{n+1}-1)$.
\end{proof}

\begin{lemma}[\texttt{gaussBinom\_mul\_qDen}]\label{lem:prod}
In every commutative ring,
$\qb{n}{k}(x)\prod_{i<k}(1-x^{i+1})=\prod_{i<k}(1-x^{n-i})$. Both sides are $0$ for $k>n$.
\end{lemma}
\begin{proof}
The proof is again an induction with the same structure. If $k\le n$, the step is
$1-x^{k+1}+x^{k+1}(1-x^{n-k})=1-x^{n+1}$. If $k>n$, the numerator contains the factor $1-x^0=0$.
\end{proof}

\begin{theorem}[(R1), \texttt{card\_subspaces\_eq\_gaussBinom}]
Let $V$ be a finite $\F$-vector space of dimension $n$. Then
$\#\{W\le V:\dim W=k\}=\qb{n}{k}(q)$.
\end{theorem}
\begin{proof}
Each linearly independent $k$-tuple $s$ in $V$ spans a $k$-dimensional subspace
$\mathrm{span}(s)$. The fibre of this map over a subspace $W$ is in bijection with the linearly
independent $k$-tuples of $W$ (\texttt{fibreEquiv}). Such a tuple spans $W$, because its length
equals $\dim W$. Mathlib's \texttt{card\_linearIndependent} counts the linearly independent
$k$-tuples in an $m$-dimensional space as $\prod_{i<k}(q^m-q^i)$ when $k\le m$. Summing over the
fibres gives
\[\#\{s\text{ lin.\ indep.}\}=N\cdot\prod_{i<k}(q^k-q^i)\qquad(\texttt{card\_li\_eq}),\]
where $N$ is the number of $k$-dimensional subspaces. Suppose $k\le n$. Then the left side is
$P_q(n,k)$. By Lemma~\ref{lem:li} we get $N\cdot P_q(k,k)=\qb{n}{k}(q)\cdot P_q(k,k)$ in
$\mathbb{Z}$. Since $q\ge2$, $P_q(k,k)\ne0$, and we may cancel. Now suppose $k>n$. Then $N=0$,
and $P_q(n,k)=0$, so $\qb{n}{k}(q)=0$ as well.
\end{proof}

(R2) follows from (R1). The first rule is the defining recursion, evaluated at $q$
(\texttt{subCount\_pascal}). For the second rule (\texttt{gaussBinom\_pascal'},
\texttt{subCount\_pascal'}), multiply both sides by $\prod_{i\le k}(1-q^{i+1})\ne0$ and use
Lemma~\ref{lem:prod}. This is the second Pascal analog in [W]. Lemma~\ref{lem:prod} also gives
the product formula for the counts (\texttt{subCount\_product\_formula}).

\begin{theorem}[(R3), \texttt{card\_subspaces\_interval}]
Let $U\le T\le V$ and $\dim U\le k$. Then
$\#\{W: U\le W\le T,\ \dim W=k\}=\qb{\dim T-\dim U}{k-\dim U}(q)$.
\end{theorem}
\begin{proof}
First the case $T=V$ (\texttt{card\_subspaces\_above}). The maps $W'\mapsto\pi^{-1}(W')$ and
$W\mapsto\pi(W)$ are inverse bijections. Here $\pi:V\to V/U$ is the quotient map, $W'$ ranges
over subspaces of $V/U$, and $W$ over subspaces containing $U$ (\texttt{aboveEquiv}). By
rank--nullity, $\dim\pi^{-1}(W')=\dim W'+\dim U$ (\texttt{finrank\_comap\_mkQ}). Apply (R1) to
$V/U$, which has dimension $\dim V-\dim U$. For general $T$, the maps given by pulling back and
pushing forward along $T\hookrightarrow V$ identify the interval $[U,T]$ in $V$ with the
subspaces of $T$ that contain $U$ (\texttt{intervalEquiv}). These maps preserve dimension.
\end{proof}

\section{Lean formalization}

The file is \texttt{lean/Conjecture2490/Basic.lean}, namespace \texttt{C2490}, 365 lines. Its
only import is \texttt{Mathlib}. The main theorem \texttt{conjecture\_2490} is a conjunction.
It states (R1) for $\F^n$ (\texttt{Fin n -> F}) for all $n,k$, both rules of (R2), and (R3) for
every finite $\F$-vector space $V$, for an arbitrary finite field $\F$. The supporting theorems
are \texttt{card\_subspaces\_eq\_gaussBinom}, \texttt{subCount\_eq},
\texttt{subCount\_pascal}, \texttt{subCount\_pascal'}, \texttt{subCount\_product\_formula},
\texttt{gaussBinom\_mul\_qDen}, \texttt{qDen\_X\_ne\_zero} and \texttt{card\_subspaces\_interval}.
As a sanity check, an \texttt{example} confirms that $N_2(4,2)=35$.

\textbf{Axiom audit.} \texttt{lake build} succeeds. \texttt{\#print axioms} reports
\texttt{[propext, Classical.choice, Quot.sound]} for \texttt{conjecture\_2490},
\texttt{card\_subspaces\_interval}, \texttt{subCount\_product\_formula},
\texttt{qDen\_X\_ne\_zero} and \texttt{gaussBinom\_mul\_qDen}. The file contains no
\texttt{sorry} and no \texttt{native\_decide}, and it adds no axioms.

\textbf{Related package.} Conjecture 00000002478 has a separate package. It covers Gauss's
q-binomial expansion $\prod_{i<n}(1+q^it)=\sum_k q^{\binom k2}\qb nk t^k$ and its counting
interpretation. This package covers the subspace count itself, both Pascal rules, and interval
closure in the subspace lattice. It does not treat the expansion.

\section*{Sources}
[W] Wikipedia, ``Gaussian binomial coefficient'',
\url{https://en.wikipedia.org/wiki/Gaussian_binomial_coefficient} (retrieved 2026-10-05). It
gives the product-formula definition and the two analogs of Pascal's identity. It also states
that $\binom nk_q$ ``counts the number of k-dimensional vector subspaces of an n-dimensional
vector space over'' $\F_q$.

Mathlib: \texttt{card\_linearIndependent}, in
\path{Mathlib/LinearAlgebra/Matrix/GeneralLinearGroup/Card.lean}.

\end{document}
